\scanpage{40}\mspage{\unsure{33|35}}
% archive sheet number: 39
chose pour le formalisme homotopique dans les catégories de fractions $\operatorname{Ho}_{W}$ etc. Mais on a \inserted{pourtant} ceci : soit
\diagram{\begin{tikzcd}[ampersand replacement=\&,row sep=0pt,column sep=small]
\&\&1\\
\unsure{\Upsilon|\Psi}={} \& 0\arrow[ur]\arrow[dr]\&\\
\&\&2
\end{tikzcd}}
, et supposons que \deleted{$(M,W)$ \uline{et}} \inserted{$(M,W)$ \uuline{et}} $(M(\Upsilon),W(\Upsilon))$ \deleted{\unsure{admette}}admettent une structure de Quillen faible exacte, \deleted{\unsure{et} $(M,W)$} Alors dans $M$ les carrés $W$-\deleted{\unsure{co}}cartésiens \unsure{sont|\deleted{étant}} \inserted{$W$-}exacts, \inserted{d'où p. \unsure{ex. ds|ex. des|ex. dans} $M$ \unsure{des} suites \uline{exactes} \inserted{longues} de cosuspension (\deleted{loop} exact loop sequences).} \deleted{Dualement, si (\unsure{Il suffit}} (\unsure{cette} hypothèse sur $(M(\Upsilon),W(\Upsilon))$ sert uniquement à assurer que toute flèche de $\operatorname{Ho}_{W(\Upsilon)}=M(\Upsilon)W(\Upsilon)^{-1}$ est isomorphe à une flèche provenant d'une flèche de $M(\Upsilon)$ lui-même. \deleted{\unsure{)}} Mais c'est là une propriété apparemment très délicate, et je ne connais d'autre façon de l'obtenir que via l'existence d'une structure de Quillen faible exacte.)
% slanted
\margin{\deleted{il semble que je me suis \unsure{embrouillé}, et je doute que \unsure{le|ce} résultat \unsure{indiqué|ici} soit exact.}}

Dualement, si $(M,W)$ et $(M(\unsure{\Psi|\Phi}),W(\Psi))$ admettent des structures de Quillen faibles \unsure{exacte|exactes}, alors les carrés $W$-cocartésiens de $M$ sont $W$-exacts, et on trouve (si $M$ est pointée) des suites exactes \uline{longues} de suspension.
